%==============================================================================
%  PLANTILLA — INFORME DE TRABAJO FINAL DE GRADUACIÓN (TFG)
%  Instituto Tecnológico de Costa Rica
%  Área Académica de Ingeniería en Computadores
%
%  Basada en la plantilla oficial Plantilla_informe_TFG.docx y en el documento
%  "Lineamientos trabajos de graduación V4.1 (CE5601)".
%
%  Compilar con:   latexmk -pdf Plantilla_Informe_TFG.tex
%  (encadena pdflatex -> biber -> pdflatex -> pdflatex, que es lo que hace
%  falta para que cuadren los índices y la bibliografía IEEE).
%
%  Estándar de citación: IEEE. Las fuentes se administran en referencias.bib.
%
%  Requiere:  ./media con image1.png (logo TEC)
%             ./referencias.bib
%==============================================================================

\documentclass[12pt,letterpaper,oneside]{article}

%------------------------------------------------------------------ Idioma ---
\usepackage[T1]{fontenc}
\usepackage[utf8]{inputenc}
\usepackage{lmodern}
% es-tabla: fuerza «Tabla» en lugar de «Cuadro» (predeterminado de babel-spanish).
% es-noquoting: evita que babel se apropie de los signos " y < >.
\usepackage[spanish,es-tabla,es-noquoting]{babel}

%----------------------------------------------------- Página y espaciado ----
\usepackage[letterpaper,margin=1in]{geometry}
\usepackage{setspace}
\onehalfspacing                 % interlineado 1,5
\usepackage{parskip}            % separación entre párrafos, sin sangría

%------------------------------------------------------- Pie de página -------
\usepackage{fancyhdr}
\pagestyle{fancy}
\fancyhf{}
\renewcommand{\headrulewidth}{0pt}
\fancyfoot[R]{\thepage}         % número de página a la derecha (como el .docx)
\fancypagestyle{plain}{\fancyhf{}\renewcommand{\headrulewidth}{0pt}\fancyfoot[R]{\thepage}}

%------------------------------------------------- Tablas, figuras, código ---
\usepackage{graphicx}
\graphicspath{{media/}{figuras/}}
\usepackage{booktabs}
\usepackage{tabularx}
\usepackage{longtable}
\usepackage{array}
\usepackage{float}
\usepackage{enumitem}
\setlist{nosep,topsep=4pt}

\usepackage{caption}
\captionsetup{font=small,labelfont=bf,labelsep=colon,skip=6pt}
\captionsetup[table]{position=above}
\captionsetup[figure]{position=below}

% Columnas de ancho elástico para tabularx:
%   C = centrada,  L = alineada a la izquierda
% Anteponga \ancho{f} para darle a una columna f veces el ancho base
% (la suma de los factores debe ser igual al número de columnas X).
\newcolumntype{C}{>{\centering\arraybackslash}X}
\newcolumntype{L}{>{\raggedright\arraybackslash}X}
\newcommand{\ancho}[1]{\hsize=#1\hsize\linewidth=\hsize}

% Lista para usar DENTRO de una celda de tabla: sin espacio superior, para que
% la primera viñeta quede alineada con el texto de las demás columnas.
\newlist{celdaitem}{itemize}{1}
\setlist[celdaitem]{label=\textbullet,leftmargin=1.2em,nosep,topsep=0pt,
                    partopsep=0pt,parsep=0pt,itemsep=2pt}

% Un poco más de aire entre las filas de las tablas
\renewcommand{\arraystretch}{1.15}

%------------------------------------------------------------- Títulos -------
\usepackage{titlesec}
\titleformat{\section}{\normalfont\Large\bfseries}{\thesection.}{0.6em}{}
\titleformat{\subsection}{\normalfont\large\bfseries}{\thesubsection.}{0.6em}{}
\titleformat{\subsubsection}{\normalfont\normalsize\bfseries}{\thesubsubsection.}{0.6em}{}
\titlespacing*{\section}{0pt}{18pt}{10pt}
\titlespacing*{\subsection}{0pt}{14pt}{8pt}
\titlespacing*{\subsubsection}{0pt}{12pt}{6pt}
\setcounter{secnumdepth}{3}     % numeración hasta 1.1.1
\setcounter{tocdepth}{3}        % índice general hasta 1.1.1

%--------------------------------------------------------- Hiperenlaces ------
\usepackage[hidelinks,unicode]{hyperref}
\hypersetup{
  pdftitle   = {Generalización de framework de selección automática de pipelines de
                preprocesamiento a múltiples tareas de aprendizaje automático mediante
                meta-aprendizaje},
  pdfauthor  = {Ricardo Antonio Borbón Mena},
  pdfsubject = {Informe de Trabajo Final de Graduación},
  pdflang    = {es-CR},
  bookmarksnumbered = true
}
\usepackage{bookmark}

%-------------------------------------------- Comando para el texto guía -----
% Todo el texto de ayuda de la plantilla va dentro de \guia{...}. Aparece en
% gris y cursiva para distinguirlo. BÓRRELO conforme escriba su informe; al
% terminar no debe quedar ningún \guia en el documento.
\usepackage{xcolor}
\newcommand{\guia}[1]{{\itshape\color{gray!75!black}#1}}

% Atajo para escribir nombres de comandos dentro del texto guía
\newcommand{\cmd}[1]{\texttt{\textbackslash #1}}

%------------------------------------------------------- Bibliografía --------
% ESTÁNDAR DE CITACIÓN: IEEE  (decidido para este TFG; los lineamientos exigen
% usar uno solo — APA 7 o IEEE — de forma rigurosa en TODO el documento).
%
% Las referencias se administran en el archivo referencias.bib.
% En el texto se cita con \cite{clave}, que produce marcas numéricas [1], [2]…
% Requiere compilar con biber (latexmk lo hace automáticamente).
\usepackage{csquotes}   % requerido por biblatex cuando se usa babel
\usepackage[backend=biber,style=ieee,sorting=none]{biblatex}
\addbibresource{referencias.bib}

% sorting=none: las referencias se numeran en el orden en que aparecen citadas
% en el texto, como manda IEEE.

%==============================================================================
\begin{document}
%==============================================================================

%------------------------------------------------------------- PORTADA -------
% La portada cuenta como página 1, pero su número no debe aparecer.
\begin{titlepage}
\thispagestyle{empty}
\centering

{\large\bfseries Instituto Tecnológico de Costa Rica\par}
\vspace{0.4cm}
{\large\bfseries Área Académica de Ingeniería en Computadores\par}
\vspace{0.4cm}
{\large\bfseries Programa de Licenciatura en Ingeniería en Computadores\par}

\vspace{1.2cm}
\includegraphics[width=7.7cm]{image1.png}

\vspace{1.5cm}
{\LARGE\bfseries Generalización de framework de selección automática de pipelines
de preprocesamiento a múltiples tareas de aprendizaje automático mediante
meta-aprendizaje\par}

\vspace{1.5cm}
{\large\bfseries Informe de Trabajo de Graduación para optar por el título de
Ingeniero en Computadores con grado académico de Licenciatura\par}

\vfill
{\large\bfseries Ricardo Antonio Borbón Mena\par}
\vspace{1cm}
{\large\bfseries Cartago, noviembre, 2026\par}

\end{titlepage}

% A partir de aquí sí se numeran las páginas; la portada fue la número 1.
\setcounter{page}{2}

%--------------------------------------------------- HOJA DE APROBACIÓN ------
\section*{Hoja de aprobación}
\addcontentsline{toc}{section}{Hoja de aprobación}

\guia{Esta sección será completada el día de la defensa.}

\guia{Debe contener las firmas del profesor asesor y del tribunal evaluador.}

\guia{El estudiante deberá digitalizarla e incluirla en el documento final.}

\guia{Debe de borrar todo este texto, incluso el título; solo debe de ir la
hoja de aprobación.}

\clearpage
%--------------------------------------------------------- DEDICATORIA -------
\section*{Dedicatoria}
\addcontentsline{toc}{section}{Dedicatoria}

\guia{En esta sección puede dedicar su trabajo a una o varias personas. Debe
ser breve y ocupar únicamente una página.}

\clearpage
%------------------------------------------------------ AGRADECIMIENTOS ------
\section*{Agradecimientos}
\addcontentsline{toc}{section}{Agradecimientos}

\guia{Esta sección es una extensión de la dedicatoria. Incluya las personas,
organizaciones o instituciones que contribuyeron al desarrollo del proyecto.
Debe redactarse de forma formal y ocupar una página.}

\clearpage
%------------------------------------------------------------- RESUMEN -------
\section*{Resumen}
\addcontentsline{toc}{section}{Resumen}

\guia{Redacte un resumen claro del proyecto en un máximo de 300 palabras. Debe
incluir el \textbf{problema abordado}, el \textbf{objetivo del proyecto}, la
\textbf{metodología utilizada}, los \textbf{resultados principales}, entre
otros puntos que considere importante. A su vez, se debe incluir 6 palabras
clave, de lo general a lo específico.}

\vspace{1em}
\noindent\textbf{Palabras clave:} \guia{palabra 1, palabra 2, palabra 3,
palabra 4, palabra 5, palabra 6.}

\clearpage
%------------------------------------------------------------ ABSTRACT -------
\section*{Abstract}
\addcontentsline{toc}{section}{Abstract}

\guia{Traducción fiel del resumen al idioma inglés. Debe mantener exactamente
el mismo contenido y estructura. Incluya también las keywords utilizadas en el
resumen.}

\vspace{1em}
\noindent\textbf{Keywords:} \guia{keyword 1, keyword 2, keyword 3, keyword 4,
keyword 5, keyword 6.}

\clearpage
%------------------------------------------------------- ÍNDICE GENERAL ------
% Se genera automáticamente a partir de los títulos. Compile dos o tres veces.
\renewcommand{\contentsname}{ÍNDICE GENERAL}
\phantomsection
\addcontentsline{toc}{section}{ÍNDICE GENERAL}
\begingroup
\singlespacing
\tableofcontents
\endgroup

\clearpage
%----------------------------------------------------- ÍNDICE DE TABLAS ------
\renewcommand{\listtablename}{ÍNDICE DE TABLAS}
\phantomsection
\addcontentsline{toc}{section}{ÍNDICE DE TABLAS}
\begingroup
\singlespacing
\listoftables
\endgroup

\clearpage
%---------------------------------------------------- ÍNDICE DE FIGURAS ------
\renewcommand{\listfigurename}{ÍNDICE DE FIGURAS}
\phantomsection
\addcontentsline{toc}{section}{ÍNDICE DE FIGURAS}
\begingroup
\singlespacing
\listoffigures
\endgroup

\clearpage
%---------------------------------------------------- SIGLAS Y ACRÓNIMOS -----
\section*{Siglas y acrónimos}
\addcontentsline{toc}{section}{Siglas y acrónimos}

\begin{description}[leftmargin=3.6cm,style=sameline,font=\bfseries]
  \item[ADR] Registro de decisión de arquitectura (\emph{Architecture Decision Record})
  \item[API] Interfaz de programación de aplicaciones (\emph{Application Programming Interface})
  \item[AutoML] Aprendizaje automático automatizado (\emph{Automated Machine Learning})
  \item[AUC-ROC] Área bajo la curva característica operativa del receptor
        (\emph{Area Under the Receiver Operating Characteristic Curve})
  \item[CC18] \emph{OpenML Curated Classification 2018}, colección curada de
        conjuntos de datos de clasificación
  \item[CSV] Valores separados por comas (\emph{Comma-Separated Values})
  \item[IA] Inteligencia artificial
  \item[IQR] Rango interquartílico (\emph{Interquartile Range})
  \item[kNN] $k$ vecinos más cercanos (\emph{$k$-Nearest Neighbors})
  \item[LODO] Exclusión de conjunto de datos (\emph{Leave-One-Dataset-Out})
  \item[MVTec AD] \emph{MVTec Anomaly Detection}, conjunto de datos de referencia
        para detección de anomalías industriales
  \item[TEC] Instituto Tecnológico de Costa Rica
  \item[TFG] Trabajo Final de Graduación
  \item[TPE] Estimador de densidad de árbol de Parzen
        (\emph{Tree-structured Parzen Estimator})
  \item[UCI] Repositorio de aprendizaje automático de la Universidad de California,
        Irvine
\end{description}

\clearpage
%=============================== CUERPO DEL INFORME ===========================

\section{Introducción}
\label{sec:introduccion}

El preprocesamiento de datos es una de las etapas que más incide en el desempeño de
los modelos de aprendizaje automático \cite{garcia2015preprocessing}. La elección de
las transformaciones aplicadas a los datos de entrada, su orden y sus parámetros
determina en gran medida la calidad de las predicciones, y suele exigir un proceso
manual de prueba y error que consume tiempo de experimentación y conocimiento experto.
La automatización de esta etapa, dentro del campo conocido como aprendizaje automático
automatizado (AutoML), busca reducir ese esfuerzo seleccionando de forma sistemática la
combinación de transformaciones que maximiza una métrica objetivo
\cite{hutter2019automl, feurer2015autosklearn}.

Este proyecto parte de una librería existente, \textit{AutoPipeAD}
\cite{borbon2026autopipead}, desarrollada por el autor en el curso CE-1114 y cuyo
propósito es automatizar la selección de pipelines de preprocesamiento para la
detección de anomalías en conjuntos de datos de imágenes. La librería define un espacio
de transformaciones componibles, genera y evalúa combinaciones mediante optimización
bayesiana con Optuna \cite{akiba2019optuna}, ejecuta las evaluaciones en paralelo y
selecciona el pipeline de mejor desempeño, generando reportes reproducibles. Está
organizada en una arquitectura modular de seis componentes desacoplados y fue probada
sobre conjuntos de datos de referencia como MVTec AD \cite{bergmann2019mvtec} y
CIFAR-10 \cite{krizhevsky2009cifar} con anomalías sintéticas.

El presente trabajo final de graduación, evoluciona esa librería hacia
\textit{AutoPipe}: un framework agnóstico a la tarea de aprendizaje, con una capa de
meta-aprendizaje que recomienda configuraciones iniciales de búsqueda a partir de las
características descriptivas del conjunto de datos utilizado.

\subsubsection{Descripción de la organización}

El proyecto se desarrolla en el Instituto Tecnológico de Costa Rica (TEC), institución
pública de educación superior universitaria creada mediante la Ley No. 4777 del 10 de
junio de 1971, dedicada a la docencia, la investigación y la extensión en los campos de
la ciencia y la tecnología. Su campus central se ubica en Cartago, Costa Rica, y cuenta
con sedes y centros académicos adicionales en distintas regiones del país.

Dentro del TEC, el proyecto se enmarca en la unidad académica responsable del programa
de Licenciatura en Ingeniería en Computadores, que articula la formación en las áreas de
electrónica digital, arquitectura de computadores, sistemas embebidos, redes, software e
inteligencia artificial. El trabajo se desarrolla bajo la supervisión de un profesor
asesor asignado por la unidad académica.

A diferencia de un proyecto realizado con una contraparte empresarial, este trabajo
tiene un carácter académico e investigativo: su contraparte es la propia unidad
académica, y su producto es una librería de software de uso abierto orientada a la
comunidad de investigadores e ingenieros de aprendizaje automático. En consecuencia, no
existe información confidencial de una organización externa que deba protegerse, y todos
los artefactos del proyecto —código, documentación y resultados experimentales— son
públicos.

\subsubsection{Descripción del área de conocimiento del proyecto}

El proyecto se ubica en el área de conocimiento de inteligencia artificial y aprendizaje
automático, dentro del eje de computación y software de la carrera de Ingeniería en
Computadores. Aborda el diseño y la implementación de un framework de software para la
selección automática de pipelines de preprocesamiento, integrando tres cuerpos de
técnicas:

\begin{itemize}
  \item \textbf{Optimización de hiperparámetros.} La búsqueda del pipeline se plantea
        como un problema de optimización sobre un espacio de configuraciones discreto,
        resuelto mediante optimización bayesiana \cite{snoek2012practical} en lugar de
        búsqueda exhaustiva.
  \item \textbf{Meta-aprendizaje.} La transferencia de conocimiento entre conjuntos de
        datos permite orientar una búsqueda nueva con la experiencia acumulada en
        búsquedas anteriores \cite{vanschoren2018metalearning, brazdil2009metalearning}.
  \item \textbf{Ingeniería de software y experimentación reproducible.} El producto es
        una librería con contratos de extensión explícitos, suite de pruebas de
        regresión y control de semillas, sobre la que se ejecuta un protocolo
        experimental definido de antemano.
\end{itemize}

La combinación de las tres es lo que caracteriza el aporte desde la Ingeniería en
Computadores: no se trata solo de aplicar un algoritmo de aprendizaje, sino de diseñar
la arquitectura de software que hace que ese algoritmo sea extensible, verificable y
reproducible.

\subsubsection{Trabajos similares encontrados}

La selección automática de pipelines de preprocesamiento es un problema estudiado dentro
del campo del AutoML \cite{hutter2019automl}. La Tabla~\ref{tab:trabajos-similares}
compara los trabajos directamente relacionados con este proyecto.

\begin{table}[H]
\centering
\small
\caption{Comparación con trabajos previos relacionados}
\label{tab:trabajos-similares}
\begin{tabularx}{\textwidth}{>{\ancho{0.55}}L>{\ancho{1.0}}L>{\ancho{1.45}}L}
\toprule
\textbf{Trabajo} & \textbf{Enfoque} & \textbf{Relación con este proyecto} \\
\midrule
\textit{auto-sklearn}
\cite{feurer2015autosklearn, feurer2015warmstart} &
Sistema AutoML que selecciona conjuntamente el modelo y su preprocesamiento dentro de un
catálogo cerrado de estimadores, e inicializa la optimización bayesiana con
configuraciones recomendadas por meta-aprendizaje. &
Es el antecedente más cercano y la fuente del mecanismo de arranque en caliente que este
proyecto adopta. Se diferencia en tres puntos: el modelo base permanece como componente
externo configurable en lugar de elegirse de un catálogo; el espacio de búsqueda se
compone de transformaciones extensibles por la persona que usa la librería, de modo que
la capa de meta-aprendizaje se especializa exclusivamente en el preprocesamiento; y se
incluye la detección de anomalías, familia de tareas que las herramientas generalistas
de AutoML no cubren. \\
\addlinespace
Inicialización de la optimización bayesiana mediante meta-aprendizaje
\cite{feurer2015warmstart} &
Demuestra que inyectar configuraciones iniciales derivadas de meta-características de
conjuntos de datos previos acelera la convergencia de la optimización bayesiana. &
Aporta el fundamento de la hipótesis central del proyecto. Este trabajo lo aplica a un
espacio de búsqueda distinto —el de pipelines de preprocesamiento— y lo valida con un
protocolo propio de exclusión de conjunto de datos. \\
\addlinespace
\textit{AutoPipeAD} \cite{borbon2026autopipead} &
Librería del propio autor que automatiza la selección de pipelines de preprocesamiento
para detección de anomalías en imágenes, con búsqueda en malla y bayesiana, evaluación
paralela y reportes reproducibles. &
Es el punto de partida directo: aporta la arquitectura modular, el motor de búsqueda y la
infraestructura de evaluación. Sus limitaciones —una sola tarea, una sola modalidad de
datos y ausencia de memoria entre ejecuciones— son precisamente lo que este proyecto
resuelve. \\
\bottomrule
\end{tabularx}
\end{table}

\subsection{Planteamiento del problema}

\subsubsection{Contexto del problema}

La construcción de un modelo de aprendizaje automático exige, antes del entrenamiento,
decidir cómo se preprocesan los datos: qué transformaciones se aplican —por ejemplo
normalización, escalado, codificación de variables categóricas, imputación de valores
faltantes, reducción de ruido o redimensionamiento—, en qué orden y con qué parámetros
\cite{garcia2015preprocessing}. El espacio de combinaciones posibles es amplio, y la
mejor combinación depende tanto de la tarea —clasificación, agrupamiento, detección de
anomalías \cite{chandola2009anomaly}— como de las características particulares del
conjunto de datos. En la práctica, esta selección se realiza de forma manual y guiada
por la experiencia, lo que la vuelve costosa en tiempo y dependiente de conocimiento
experto.

La librería \textit{AutoPipeAD} automatiza esa selección, pero únicamente para la tarea
de detección de anomalías sobre conjuntos de datos de imágenes. Su motor de búsqueda, su
mecanismo de paralelización, su abstracción de pipelines componibles y su generación de
reportes reproducibles son, por diseño, independientes del dominio. El acoplamiento a
imágenes y a la detección de anomalías se concentra, en cambio, en unos pocos puntos
concretos del código, lo que impide que la herramienta cumpla su propósito original de
ser de propósito general.

Por otra parte, cada vez que la librería se enfrenta a un conjunto de datos nuevo
ejecuta la búsqueda partiendo desde cero, sin aprovechar información de búsquedas
anteriores sobre conjuntos de datos similares. Dado que la búsqueda automática consume
tiempo de cómputo proporcional al número de combinaciones evaluadas, repetir el proceso
completo para cada conjunto de datos nuevo representa una oportunidad de mejora
identificable.

\subsubsection{Justificación del problema}

\textbf{Innovación.} La generalización del framework a múltiples tareas, combinada con
una capa de meta-aprendizaje que recomienda configuraciones iniciales de búsqueda a
partir de meta-características del conjunto de datos, constituye un aporte que va más
allá de la ingeniería incremental. Mientras que la selección automática de pipelines es
un problema estudiado en AutoML \cite{hutter2019automl}, la transferencia de conocimiento
entre conjuntos de datos para acelerar específicamente la selección de pipelines de
preprocesamiento sigue siendo un área con margen de aporte
\cite{brazdil2009metalearning}, y permite formular una hipótesis verificable.

\textbf{Evidencia del problema.} El acoplamiento al dominio de imágenes y a la detección
de anomalías es directamente observable en el diseño de la librería: su componente de
carga de datos asume una estructura de directorios propia de imágenes y su componente de
evaluación calcula métricas específicas de detección de anomalías. El arranque en frío es
igualmente verificable: el motor de búsqueda no consume ni produce ninguna base de
conocimiento entre ejecuciones, por lo que cada nueva búsqueda parte siempre desde cero.
La auditoría de acoplamiento realizada como parte de este trabajo, descrita en la
Sección~\ref{sec:trabajo-realizado}, confirma y precisa ese diagnóstico.

\textbf{Valor generado.} Para la comunidad de investigadores e ingenieros de aprendizaje
automático, una herramienta que automatice el preprocesamiento de forma agnóstica a la
tarea reduce el tiempo de experimentación y la dependencia de conocimiento experto. La
capa de meta-aprendizaje añade valor al reducir el costo de cómputo de la búsqueda, lo
que se traduce en menor tiempo y menor consumo energético por experimento. Al área de
conocimiento aporta un artefacto de software reutilizable y un estudio experimental
reproducible sobre la efectividad del arranque en caliente meta-aprendido, sea el
resultado positivo o negativo.

\subsubsection{Enunciado del problema}

¿Cómo generalizar la selección automática de pipelines de preprocesamiento a múltiples
tareas de aprendizaje automático y reducir el costo de la búsqueda en conjuntos de datos
nuevos aprovechando, mediante meta-aprendizaje, el conocimiento acumulado sobre conjuntos
de datos previos?

\subsection{Objetivos del proyecto}

\subsubsection{Objetivo general}

Generalizar el framework AutoPipeAD de selección automática de pipelines de
preprocesamiento a múltiples tareas de aprendizaje automático incorporando una capa de
meta-aprendizaje que recomiende configuraciones iniciales de búsqueda a partir de las
meta-características del conjunto de datos, reduciendo así el costo de exploración
en conjuntos de datos nuevos manteniendo la reproducibilidad de los resultados.

\subsubsection{Objetivos específicos}

\begin{enumerate}
  \item Disponer de un núcleo agnóstico a la tarea que resuelva detección de anomalías y
        clasificación tabular sobre una misma arquitectura, verificado por la suite de
        pruebas de regresión de la librería base y por pruebas que demuestren la
        ausencia de fuga de información entre los subconjuntos de entrenamiento y
        validación.

  \item Contar con una base de meta-conocimiento que asocie, para un corpus de conjuntos
        de datos públicos, sus meta-características con los pipelines de mejor desempeño
        hallados por la búsqueda, completa y reproducible bajo semillas controladas.

  \item Obtener un recomendador meta-aprendido capaz de producir configuraciones
        iniciales válidas y ejecutables para un conjunto de datos no visto, integradas en
        el motor de búsqueda como ensayos iniciales.

  \item Determinar, con evidencia estadística, si el arranque en caliente meta-aprendido
        reduce el costo de la búsqueda frente al arranque en frío, mediante validación
        cruzada por exclusión de conjunto de datos y reportes reproducibles.
\end{enumerate}

\subsection{Alcances, entregables y limitaciones del proyecto}
\label{sec:alcances}

\subsubsection{Alcance}

Para administrar el presupuesto de 336 horas con un margen de riesgo explícito, el
alcance se organiza en dos niveles.

El \textbf{alcance mínimo comprometido} constituye el criterio de completitud del
proyecto: el núcleo agnóstico a la tarea con soporte para detección de anomalías —caso de
continuidad de \textit{AutoPipeAD}— y clasificación sobre datos tabulares; el extractor
de meta-características; la base de meta-conocimiento sobre un corpus acotado de
conjuntos de datos públicos; el recomendador de arranque en caliente integrado con el
motor de búsqueda; la validación experimental por exclusión de conjunto de datos; y la
documentación técnica con el manual de uso.

Las \textbf{extensiones deseables} se abordan únicamente si el avance verificado del
alcance mínimo lo permite: el soporte de agrupamiento, cuya evaluación no supervisada
requiere un protocolo propio; la ampliación del corpus de meta-conocimiento; y un
meta-modelo de recomendación más sofisticado que el recomendador por vecinos más
cercanos. La decisión de incorporarlas o descartarlas se toma en un punto de control de
alcance situado en la séptima semana del cronograma.

\subsubsection{Entregables}

\begin{table}[H]
\centering
\small
\caption{Entregables del proyecto}
\label{tab:entregables}
\begin{tabularx}{\textwidth}{>{\ancho{0.45}}C>{\ancho{1.0}}L>{\ancho{1.55}}L}
\toprule
\textbf{\#} & \textbf{Entregable} & \textbf{Descripción} \\
\midrule
RE1 & Núcleo agnóstico a la tarea &
Refactorización de \textit{AutoPipeAD} con una interfaz de tarea y un registro de
métricas que soportan detección de anomalías y clasificación. \\
\addlinespace
RE2 & Soporte de datos tabulares &
Cargador de conjuntos de datos tabulares y transformaciones de preprocesamiento tabular
—escalado, codificación, imputación y selección de características— construidas sobre la
extensión del contrato de transformación con semántica de ajuste y transformación. \\
\addlinespace
RE3 & Extractor de meta-características &
Componente que calcula un vector de meta-características descriptivas de un conjunto de
datos: estadísticas, teórico-informativas y de complejidad. \\
\addlinespace
RE4 & Base de meta-conocimiento &
Estructura de datos que asocia, para un corpus de conjuntos de datos públicos, sus
meta-características con los pipelines de mejor desempeño hallados por la búsqueda. \\
\addlinespace
RE5 & Recomendador meta-aprendido &
Modelo que, dadas las meta-características de un conjunto de datos nuevo, produce
configuraciones iniciales de arranque en caliente para el motor de búsqueda. \\
\addlinespace
RE6 & Experimentos y reportes comparativos &
Conjunto de experimentos reproducibles y reportes que comparan arranque en frío contra
arranque en caliente y documentan la generalización por tarea. \\
\addlinespace
RE7 & Documentación y código entregado &
Documentación técnica, manual de uso y código fuente del framework en repositorio. \\
\bottomrule
\end{tabularx}
\end{table}

\subsubsection{Limitaciones}

\begin{itemize}
  \item El alcance comprometido se restringe a datos tabulares y a las familias de tareas
        de clasificación y detección de anomalías. No se cubren texto, series temporales
        ni grafos.
  \item El framework no entrena ni optimiza el modelo base, que permanece como componente
        externo configurable a través de una interfaz compatible con scikit-learn
        \cite{sklearn-docs}. El único modelo entrenado dentro del proyecto es el
        meta-modelo interno del recomendador.
  \item El aumento de datos mediante técnicas generativas queda fuera de alcance por las
        restricciones de tiempo y de hardware, y se documenta como trabajo futuro.
  \item La paralelización permanece limitada al hardware local; no se contempla
        distribución en clústeres ni en infraestructura de nube.
  \item El tamaño del corpus de meta-conocimiento está acotado por el presupuesto de
        cómputo y de tiempo, lo que limita el alcance estadístico de las conclusiones
        sobre el meta-aprendizaje.
  \item La validación se realiza sobre conjuntos de datos públicos; no se garantiza la
        transferencia de los resultados a conjuntos propietarios o a dominios
        sustancialmente distintos de los del corpus.
  \item El proyecto se ejecuta de forma individual con un presupuesto de 336 horas
        distribuidas en doce semanas, lo que constituye la restricción más fuerte sobre
        la profundidad alcanzable en cada componente.
\end{itemize}

\clearpage
%------------------------------------------------------------------------------
\section{Marco de referencia teórico}
\label{sec:marco-teorico}

Este capítulo establece el sistema de conceptos sobre el que se apoya el proyecto. Se
organiza siguiendo la cadena lógica que va del problema a la solución propuesta: primero
el preprocesamiento como etapa determinante del desempeño; luego su automatización dentro
del AutoML y el mecanismo de búsqueda empleado; después el meta-aprendizaje como forma de
transferir conocimiento entre conjuntos de datos; y finalmente los dos elementos que
permiten medir el resultado, que son las métricas por tarea y el protocolo de validación.

La Figura~\ref{fig:mapa-conceptual} presenta el mapa conceptual del capítulo.

\begin{figure}[H]
\centering
\begin{tabularx}{0.95\textwidth}{C}
\toprule
\textbf{Preprocesamiento de datos} \\
\emph{Determina el desempeño del modelo; su selección manual es costosa} \\
\midrule
$\downarrow$ \\
\midrule
\textbf{AutoML: selección automática de pipelines} \\
\emph{Búsqueda sobre un espacio de transformaciones componibles} \\
\midrule
$\downarrow$ \\
\midrule
\textbf{Optimización bayesiana (TPE)} \\
\emph{Explora el espacio con menos evaluaciones que la búsqueda exhaustiva} \\
\midrule
$\downarrow$ \\
\midrule
\textbf{Meta-aprendizaje y arranque en caliente} \\
\emph{Meta-características $\rightarrow$ recomendación $\rightarrow$ ensayos iniciales} \\
\midrule
$\downarrow$ \\
\midrule
\textbf{Medición: métricas por tarea y validación por exclusión de conjunto de datos} \\
\emph{Permite decidir si el arranque en caliente reduce el costo de la búsqueda} \\
\bottomrule
\end{tabularx}
\caption{Mapa conceptual del marco teórico}
\label{fig:mapa-conceptual}
\end{figure}

\subsection{Preprocesamiento de datos y su efecto en el desempeño}

El preprocesamiento comprende el conjunto de operaciones que transforman los datos crudos
en una representación apta para el entrenamiento: escalado de variables numéricas,
imputación de valores faltantes, codificación de variables categóricas, selección de
características y, en el caso de imágenes, operaciones como reducción de ruido, ajuste de
contraste o redimensionamiento \cite{garcia2015preprocessing}. Su relevancia no es
accesoria: la calidad de la representación acota el desempeño alcanzable por cualquier
modelo posterior, de modo que dos pipelines distintos sobre los mismos datos y el mismo
modelo pueden producir resultados sustancialmente diferentes.

Dos propiedades del preprocesamiento resultan determinantes para el diseño de este
proyecto. La primera es que el efecto de una transformación \textbf{depende del modelo}: los
modelos basados en distancias, como $k$ vecinos más cercanos, son sensibles al escalado,
mientras que los basados en particiones del espacio de características lo son mucho menos.
La segunda es que muchas transformaciones tienen \textbf{estado}: el escalado, la imputación y
la codificación requieren estimar cantidades a partir de los datos —medias, medianas,
catálogos de categorías— antes de poder aplicarse. Esta segunda propiedad tiene una
consecuencia metodológica que se desarrolla en la Sección~\ref{sec:trabajo-realizado}: si
esas cantidades se estiman sobre el conjunto completo en lugar de solo sobre el
subconjunto de entrenamiento, se produce fuga de información y la evaluación deja de ser
válida.

\subsection{Aprendizaje automático automatizado}

El AutoML busca automatizar las decisiones de diseño que tradicionalmente toma una
persona experta: selección de modelo, ajuste de hiperparámetros y, en el caso que
interesa aquí, selección del preprocesamiento \cite{hutter2019automl}. Formalmente, la
selección de un pipeline puede plantearse como un problema de optimización sobre un
espacio de configuraciones: dado un espacio $\mathcal{P}$ de pipelines posibles, un
modelo base fijo y una métrica objetivo $m$, se busca el pipeline $p^\ast \in \mathcal{P}$
que maximiza $m$ evaluada sobre un subconjunto de validación.

El tamaño de $\mathcal{P}$ crece de forma multiplicativa con el número de pasos y de
alternativas por paso, lo que vuelve impracticable la evaluación exhaustiva en cuanto el
espacio contiene más de unos pocos pasos. De ahí la necesidad de una estrategia de
búsqueda que no evalúe todas las combinaciones.

\textit{auto-sklearn} \cite{feurer2015autosklearn} es el sistema de referencia en este
campo: combina selección de modelo y de preprocesamiento mediante optimización bayesiana,
y utiliza meta-aprendizaje para inicializar la búsqueda. Su diferencia con este proyecto,
detallada en la Tabla~\ref{tab:trabajos-similares}, es que el modelo base forma parte del
espacio de búsqueda, mientras que aquí permanece como componente externo fijo, lo que
permite atribuir las diferencias de desempeño observadas al preprocesamiento y no al
modelo.

\subsection{Optimización bayesiana como estrategia de búsqueda}

La optimización bayesiana aborda la optimización de funciones costosas de evaluar
construyendo un modelo probabilístico de la función objetivo y usándolo para decidir qué
configuración evaluar a continuación \cite{snoek2012practical}. La decisión equilibra
explotación —evaluar donde el modelo predice buenos valores— y exploración —evaluar donde
la incertidumbre es alta—, de modo que el número de evaluaciones necesarias es
típicamente mucho menor que en búsqueda aleatoria o exhaustiva.

Este proyecto emplea la implementación de Optuna \cite{akiba2019optuna}, cuyo muestreador
predeterminado es el estimador de densidad de árbol de Parzen (TPE). A diferencia de los
métodos basados en procesos gaussianos, TPE modela las densidades condicionales de las
configuraciones que produjeron buenos y malos resultados, lo que lo hace especialmente
adecuado para espacios de búsqueda con variables categóricas y condicionales, como es el
caso de un espacio de pipelines donde cada paso puede además omitirse.

Dos aspectos del comportamiento de TPE resultan relevantes para el diseño experimental.
El primero es que el muestreador no comienza a modelar de inmediato: durante un número
configurable de ensayos iniciales muestrea al azar, y solo después construye el modelo de
densidades \cite{optuna-docs}. El segundo es que Optuna permite inyectar configuraciones
concretas para que se evalúen como primeros ensayos del estudio, mecanismo sobre el que se
implementa el arranque en caliente. Ambos aspectos fueron verificados experimentalmente en
este trabajo, como se describe en la Sección~\ref{sec:trabajo-realizado}.

\subsection{Meta-aprendizaje y arranque en caliente}

El meta-aprendizaje estudia cómo aprender de la experiencia acumulada en múltiples tareas
de aprendizaje para mejorar el desempeño en tareas nuevas
\cite{vanschoren2018metalearning, brazdil2009metalearning}. En el contexto de la
optimización de hiperparámetros, la forma más directa de aplicarlo es el
\textit{arranque en caliente}: en lugar de iniciar la búsqueda desde configuraciones
aleatorias, se inicia desde configuraciones que resultaron buenas en conjuntos de datos
similares \cite{feurer2015warmstart}.

El mecanismo requiere tres piezas:

\begin{enumerate}
  \item Una \textbf{descripción cuantitativa} de cada conjunto de datos, denominada vector
        de meta-características. Suele agrupar cuatro tipos de medidas: generales, como el
        número de instancias y de atributos; estadísticas, como los momentos de las
        distribuciones marginales; teórico-informativas, como entropías e información
        mutua; y de complejidad, que son medidas geométricas sobre la separabilidad de las
        clases.
  \item Una \textbf{base de meta-conocimiento} que asocie, para un corpus de conjuntos de
        datos, su vector de meta-características con las configuraciones que mejor
        funcionaron en ellos.
  \item Un \textbf{mecanismo de recomendación} que, dado el vector de un conjunto de datos
        nuevo, seleccione configuraciones de la base. El enfoque habitual, y el adoptado
        en este proyecto, es la búsqueda de los vecinos más cercanos en el espacio de
        meta-características.
\end{enumerate}

La extracción de meta-características se realiza con \textit{pymfe}
\cite{pymfe-docs}, librería que implementa los grupos de medidas descritos arriba. La
elección de una librería estándar en lugar de una implementación propia responde a la
mitigación de un riesgo identificado en el anteproyecto: la heterogeneidad de los
conjuntos de datos públicos dificulta la extracción uniforme de meta-características, y
una implementación establecida cubre más casos límite que una escrita para el proyecto.

\subsection{Evaluación: métricas por tarea y protocolo de validación}

Cada familia de tareas requiere métricas propias. Para la \textbf{detección de anomalías}
\cite{chandola2009anomaly}, planteada como clasificación binaria con clases muy
desbalanceadas, la métrica primaria de este proyecto es el área bajo la curva
característica operativa del receptor (AUC-ROC), que evalúa la capacidad de ordenamiento
del modelo con independencia del umbral de decisión. Para la \textbf{clasificación}, la
métrica primaria es el F1 macro, que promedia el F1 de cada clase con igual peso; a
diferencia de la exactitud, penaliza a un modelo que acierte solo en la clase
mayoritaria, propiedad necesaria cuando el corpus contiene conjuntos de datos
desbalanceados.

Para decidir si el arranque en caliente aporta valor no basta comparar el desempeño final,
porque ambas condiciones convergen a valores parecidos con presupuesto suficiente. Lo que
debe medirse es el \textbf{costo de la búsqueda}: cuántos ensayos hacen falta para
alcanzar un desempeño objetivo. La comparación se realiza mediante validación cruzada por
exclusión de conjunto de datos (LODO), que para cada conjunto de datos del corpus
construye la base de meta-conocimiento con los restantes y evalúa sobre el excluido. Este
protocolo es el análogo, a nivel de conjuntos de datos, de la validación cruzada
convencional a nivel de instancias, y su supuesto es que los conjuntos de datos del corpus
son independientes entre sí. La verificación de ese supuesto resultó ser un hallazgo
relevante de este trabajo y se discute en la Sección~\ref{sec:trabajo-realizado}.

El corpus se construye a partir de conjuntos de datos públicos de OpenML, y en particular
de la colección curada OpenML-CC18 \cite{openml-cc18}, cuya justificación como universo de
candidatos se detalla en el capítulo siguiente.

\clearpage
%------------------------------------------------------------------------------
\section{Marco metodológico}
\label{sec:marco-metodologico}

Este capítulo describe la estrategia seguida para resolver el problema y las tácticas
concretas empleadas en cada etapa. Se presenta primero la naturaleza del trabajo y la
justificación de la estrategia elegida, luego la secuencia del proceso, y finalmente el
desglose que asocia cada objetivo específico con sus actividades, entregables,
herramientas y estrategias de verificación.

\subsection{Naturaleza del trabajo y estrategia elegida}

El proyecto combina dos modalidades. Por una parte corresponde al desarrollo de un nuevo
producto, en tanto evoluciona una librería existente hacia un framework agnóstico a la
tarea e incorpora una capa de meta-aprendizaje que la herramienta actual no posee. Por
otra parte incorpora una componente de investigación teórico-práctica, dado que formula y
evalúa empíricamente una hipótesis verificable: si el arranque en caliente meta-aprendido
reduce el costo de la búsqueda frente al arranque en frío.

La estrategia adoptada responde a tres decisiones metodológicas.

\textbf{Evolución sobre reinvención.} Se reutiliza la arquitectura modular existente y se
construye sobre sus decisiones de diseño consolidadas —motor de búsqueda, paralelización y
modelo base externo configurable— en lugar de reescribir el framework. La consecuencia
práctica es que la suite de pruebas de la librería base, con 134 pruebas y una cobertura
del 97\,\%, funciona como red de regresión durante toda la refactorización: cualquier
cambio que rompa el caso de uso original se detecta de inmediato.

\textbf{Desarrollo iterativo e incremental con verificación continua.} El trabajo se
organiza en etapas secuenciales donde cada una construye sobre los resultados de la
anterior, y cada incremento se acompaña de sus pruebas. La verificación no se posterga a
una fase de validación final: la suite completa, el analizador de estilo y el verificador
de tipos se ejecutan en cada incremento.

\textbf{Decisiones de diseño y de protocolo documentadas antes de implementar o medir.}
Las decisiones arquitectónicas se registran en documentos de decisión de arquitectura
(ADR) que consignan el contexto, las alternativas consideradas, la decisión tomada y sus
consecuencias, antes de escribir el código correspondiente. De forma análoga, y por
tratarse de un trabajo con componente experimental, las decisiones del protocolo de
medición se fijan en un documento de pre-registro versionado y fechado, cerrado antes de
ejecutar los experimentos que gobierna. El propósito es impedir que las reglas de análisis
se ajusten después de conocer los resultados.

Complementariamente, antes de comprometer trabajo sobre supuestos técnicos no verificados
se ejecutan \textit{spikes}: experimentos acotados cuyo único objetivo es medir la
viabilidad de un mecanismo. Esta táctica se aplicó a los dos supuestos de mayor riesgo del
proyecto, la inyección de ensayos iniciales en el motor de búsqueda y el costo de la
extracción de meta-características.

\subsection{Secuencia del proceso}

El proceso se organiza en las ocho etapas de la
Tabla~\ref{tab:etapas}, donde cada una produce los insumos de la siguiente.

\begin{table}[H]
\centering
\small
\caption{Etapas del proceso y su encadenamiento}
\label{tab:etapas}
\begin{tabularx}{\textwidth}{>{\ancho{0.42}}C>{\ancho{1.05}}L>{\ancho{1.53}}L}
\toprule
\textbf{\#} & \textbf{Etapa} & \textbf{Insumo que produce para la siguiente} \\
\midrule
E1 & Análisis y diseño de la generalización &
Contrato de transformación con semántica de ajuste, interfaz de tarea y criterios de
inclusión del corpus. \\
\addlinespace
E2 & Implementación del núcleo agnóstico y soporte tabular &
Núcleo capaz de evaluar clasificación tabular sin fuga de información. \\
\addlinespace
E3 & Implementación del extractor de meta-características &
Vector de meta-características reproducible para cualquier conjunto de datos. \\
\addlinespace
E4 & Construcción de la base de meta-conocimiento &
Corpus descrito y sus pipelines ganadores. \\
\addlinespace
E5 & Implementación del recomendador meta-aprendido &
Configuraciones iniciales inyectables en el motor de búsqueda. \\
\addlinespace
E6 & Diseño experimental y validación &
Mediciones de costo de búsqueda en ambas condiciones. \\
\addlinespace
E7 & Generación de reportes y análisis de resultados &
Evidencia sobre la hipótesis y sobre la generalización por tarea. \\
\addlinespace
E8 & Documentación y entrega &
Documentación técnica, manual de uso y código en repositorio. \\
\bottomrule
\end{tabularx}
\end{table}

La Figura~\ref{fig:flujo-solucion} resume el flujo de la solución en ejecución, una vez
construidos todos los componentes. Los bloques sombreados corresponden a componentes
nuevos o extendidos respecto de la librería base.

\begin{figure}[H]
\centering
\begin{tabularx}{0.95\textwidth}{C}
\toprule
Conjunto de datos (tabular o imágenes) $+$ tarea de aprendizaje \\
\midrule
$\downarrow$ \\
\midrule
\textbf{Extractor de meta-características} $\rightarrow$
\textbf{recomendador meta-aprendido} $\rightarrow$ configuraciones iniciales \\
\midrule
$\downarrow$ \\
\midrule
Motor de búsqueda con arranque en caliente $\rightarrow$ evaluación paralela
$\rightarrow$ selección \\
\midrule
$\downarrow$ \\
\midrule
Pipeline de mejor desempeño $+$ reporte reproducible \\
\bottomrule
\end{tabularx}
\caption{Flujo de la solución con la capa de meta-aprendizaje}
\label{fig:flujo-solucion}
\end{figure}

\subsection{Desglose por objetivo específico}

Las Tablas~\ref{tab:desglose} y~\ref{tab:desglose-2} desglosan cada objetivo específico en
sus actividades principales, los entregables que produce, las técnicas y herramientas
empleadas y las estrategias de verificación y validación aplicadas.

\begin{table}[H]
\centering
\small
\caption{Desglose del marco metodológico: objetivos específicos 1 y 2}
\label{tab:desglose}
\begin{tabularx}{\textwidth}{>{\ancho{0.8}}L>{\ancho{1.32}}L>{\ancho{0.62}}C>{\ancho{0.9}}L>{\ancho{1.36}}L}
\toprule
\textbf{Objetivo} & \textbf{Actividades principales} & \textbf{Entrega\-bles} &
\textbf{Técnicas y herramientas} & \textbf{Verificación y validación} \\
\midrule
\textbf{OE1.} Núcleo agnóstico a la tarea con clasificación tabular &
\begin{celdaitem}
  \item Auditoría del acoplamiento a imágenes y anomalías
  \item Diseño del contrato de ajuste y transformación
  \item Diseño de la interfaz de tarea y del registro de métricas
  \item Refactorización del protocolo de evaluación
  \item Cargador y transformaciones tabulares
\end{celdaitem} &
RE1, RE2 &
Python, scikit-learn, pandas, patrón estrategia, ADR, \texttt{pytest}, \texttt{ruff},
\texttt{mypy} &
\begin{celdaitem}
  \item Suite de regresión de la librería base en verde, sin editar
  \item Prueba de ausencia de fuga entre entrenamiento y validación
  \item Prueba de aislamiento entre candidatos de la búsqueda
  \item Cobertura mínima del 70\,\% en módulos nuevos
\end{celdaitem} \\
\addlinespace
\textbf{OE2.} Base de meta-conocimiento completa y reproducible &
\begin{celdaitem}
  \item Implementación del extractor de meta-características
  \item Selección del corpus según criterios de inclusión
  \item Ejecución paralela de la búsqueda sobre el corpus
  \item Registro de meta-características y pipelines ganadores
\end{celdaitem} &
RE3, RE4 &
\textit{pymfe}, OpenML, \texttt{joblib}, Optuna &
\begin{celdaitem}
  \item Consistencia y determinismo del vector sobre conjuntos de control
  \item Completitud de la base para todo el corpus
  \item Reproducibilidad: misma semilla, mismo pipeline ganador
\end{celdaitem} \\
\bottomrule
\end{tabularx}
\end{table}

\begin{table}[H]
\centering
\small
\caption{Desglose del marco metodológico: objetivos específicos 3 y 4}
\label{tab:desglose-2}
\begin{tabularx}{\textwidth}{>{\ancho{0.8}}L>{\ancho{1.32}}L>{\ancho{0.62}}C>{\ancho{0.9}}L>{\ancho{1.36}}L}
\toprule
\textbf{Objetivo} & \textbf{Actividades principales} & \textbf{Entrega\-bles} &
\textbf{Técnicas y herramientas} & \textbf{Verificación y validación} \\
\midrule
\textbf{OE3.} Recomendador que produce configuraciones válidas y ejecutables &
\begin{celdaitem}
  \item Construcción del meta-modelo por vecinos más cercanos
  \item Generación de configuraciones a partir de meta-características
  \item Integración con el motor de búsqueda como ensayos iniciales
\end{celdaitem} &
RE5 &
scikit-learn, Optuna &
\begin{celdaitem}
  \item Validación de que cada configuración recomendada es válida en la versión vigente
        del espacio de búsqueda
  \item Evaluación preliminar sobre conjuntos de datos retenidos
\end{celdaitem} \\
\addlinespace
\textbf{OE4.} Evidencia estadística sobre la reducción del costo de búsqueda &
\begin{celdaitem}
  \item Pre-registro del protocolo antes de ejecutar
  \item Instrumentación de la validación por exclusión de conjunto de datos
  \item Ejecución de arranque en frío contra arranque en caliente
  \item Análisis estadístico y generación de reportes
\end{celdaitem} &
RE6, RE7 &
Validación LODO, prueba de rangos con signo de Wilcoxon, \texttt{scipy}, reportes
HTML y PDF &
\begin{celdaitem}
  \item Comparación pareada por conjunto de datos con significancia estadística
  \item Análisis de robustez frente a la regla de censura
  \item Verificación de la generalización por tarea
  \item Revisión de resultados con el profesor asesor
\end{celdaitem} \\
\bottomrule
\end{tabularx}
\end{table}

\subsection{Criterios de éxito}

Para que la validación sea verificable, los criterios de éxito quedaron fijados de
antemano en el anteproyecto \cite{borbon2026anteproyecto} y se resumen a continuación.

\textbf{Corpus y presupuesto de búsqueda.} Entre 15 y 25 conjuntos de datos públicos de
clasificación tabular, con licencia de uso abierto, entre 500 y 50\,000 instancias, no más
de 100 características, ausencia de datos de imagen o texto embebidos y un máximo de 10
clases. La construcción de la base de meta-conocimiento ejecuta 50 ensayos de búsqueda por
conjunto de datos, con semillas controladas y resultados cacheados.

\textbf{Métricas primarias.} F1 macro para clasificación y AUC-ROC para detección de
anomalías; las demás métricas se reportan con fines descriptivos.

\textbf{Protocolo de comparación.} Para cada conjunto de datos del corpus, la base de
meta-conocimiento se construye con los restantes y se comparan dos condiciones sobre el
excluido, ambas con un presupuesto de 30 ensayos y las mismas semillas: arranque en frío,
que es la búsqueda estándar, y arranque en caliente, la misma búsqueda con cinco
configuraciones recomendadas inyectadas como ensayos iniciales. En ambas se emplea el mismo
modelo base fijo, un clasificador sensible al preprocesamiento, de modo que las diferencias
observadas sean atribuibles al pipeline y no al modelo. Se miden dos variables: los ensayos
necesarios para alcanzar al menos el 95\,\% del mejor valor hallado por el arranque en frío
con el presupuesto completo, y el mejor valor alcanzado tras la mitad del presupuesto.

\textbf{Criterio de la hipótesis.} Se considera confirmada si el arranque en caliente
reduce la mediana de ensayos hasta el objetivo en al menos un 30\,\% respecto del arranque
en frío, con significancia estadística según la prueba de rangos con signo de Wilcoxon
pareada por conjunto de datos con $\alpha = 0{,}05$. Si el criterio no se alcanza, el
análisis riguroso de las causas y la generalización del framework constituyen por sí mismos
resultados defendibles.

\textbf{Criterios de éxito de ingeniería.} Con independencia del resultado de la hipótesis:
la suite de regresión de la librería base se mantiene en verde tras la refactorización; los
módulos nuevos alcanzan al menos un 70\,\% de cobertura de pruebas; los resultados son
reproducibles, en el sentido de que la misma semilla produce el mismo pipeline ganador; y
el flujo completo para la persona que usa la librería se mantiene expresable en no más de
20 líneas de código.

\subsection{Decisiones de protocolo precisadas durante el desarrollo}

Al implementar el protocolo se identificaron puntos que los criterios de éxito no
especificaban y que, de resolverse después de observar los resultados, habrían constituido
grados de libertad indebidos. Se fijaron por escrito en el pre-registro, sin modificar
ningún criterio ya comprometido. La Tabla~\ref{tab:decisiones-protocolo} los resume.

\begin{table}[htbp]
\centering
\small
\caption{Decisiones de protocolo precisadas y su justificación}
\label{tab:decisiones-protocolo}
\begin{tabularx}{\textwidth}{>{\ancho{0.95}}L>{\ancho{1.0}}L>{\ancho{1.05}}L}
\toprule
\textbf{Punto no especificado} & \textbf{Decisión} & \textbf{Justificación} \\
\midrule
Modelo base concreto, entre las dos opciones admitidas &
$k$ vecinos más cercanos con configuración congelada &
Opera sobre distancias, de modo que escalado, imputación y selección de características
afectan la métrica de forma sustantiva. Un modelo insensible al preprocesamiento
produciría un efecto piso que vaciaría la comparación. \\
\addlinespace
Parametrización de las transformaciones &
Toda cantidad dependiente del conjunto de datos se expresa en términos relativos &
El arranque en caliente transfiere configuraciones entre conjuntos de datos; un parámetro
absoluto no significa lo mismo en dos conjuntos distintos y puede ser inválido en el
receptor. \\
\addlinespace
Número de semillas por condición &
Tres semillas, idénticas en ambas condiciones, agregadas por mediana &
Una sola corrida produce una observación pareada ruidosa, y una conclusión que dependa de
la semilla obtenida no es defendible. No altera el presupuesto de ensayos ni el
emparejamiento. \\
\addlinespace
Casos en que una condición no alcanza el objetivo &
Se asigna el presupuesto más uno, de forma simétrica, y se reporta el recuento; se
recalcula el análisis con castigos más severos &
El arranque en frío no puede quedar censurado, porque el objetivo se define a partir de su
propio mejor valor. La regla solo puede afectar al arranque en caliente, y por eso
requiere una garantía explícita de robustez. \\
\addlinespace
Lateralidad de la prueba estadística &
Una cola en la dirección de la hipótesis, reportando también el valor de dos colas &
La hipótesis es direccional. Declarar la lateralidad de antemano y reportar ambos valores
impide que la elección se convierta en un grado de libertad oculto. \\
\bottomrule
\end{tabularx}
\end{table}

\clearpage
%------------------------------------------------------------------------------
\section{Descripción del trabajo realizado}
\label{sec:trabajo-realizado}

\subsection{Descripción del proceso de solución}

Esta sección describe el trabajo ejecutado hasta la fecha de redacción. Cubre las etapas
E1 y E2 del marco metodológico, correspondientes a los entregables RE1 y RE2, y el inicio
de RE4 y RE7. Las etapas E3 a E7 están en curso y su descripción se completará conforme
avancen.

\subsubsection{Auditoría del acoplamiento: punto de partida verificado}

El primer paso fue determinar, sobre el código y no sobre supuestos, qué partes de la
librería base dependían de imágenes y de la detección de anomalías. El anteproyecto
sostenía que el acoplamiento se concentraba en la carga de datos y en el cálculo de
métricas. La auditoría confirmó ese diagnóstico en lo esencial, pero lo corrigió en dos
direcciones.

Por una parte, redujo el trabajo previsto: el componente de partición de datos resultó
completamente independiente de la tarea, y el generador de reportes deriva las columnas de
métricas de los resultados en lugar de tomarlas de una lista fija, de modo que ninguno
requería modificación.

Por otra parte —y este fue el hallazgo determinante— identificó \textbf{dos funciones
adicionales, dentro del motor de búsqueda, específicas de la detección de anomalías y de
fallo silencioso}. La normalización de predicciones binariza la salida del modelo, de modo
que ante un problema de cinco clases convierte el vector $[0,1,2,3,4]$ en $[0,1,1,1,1]$; y
la obtención de puntajes continuos toma una sola columna de la matriz de probabilidades,
descartando cuatro de cinco columnas en un problema multiclase. Ninguna de las dos levanta
excepción: producen métricas plausibles y equivocadas. Ambos comportamientos se verificaron
ejecutando el código, no leyéndolo.

La consecuencia de diseño fue directa: la abstracción de tarea debía agrupar cuatro
responsabilidades y no dos. Separar las métricas de la normalización de predicciones
habría permitido que una tarea nueva declarara sus métricas y heredara en silencio la
normalización binaria, que es exactamente el modo de fallo documentado.

\subsubsection{Riesgo técnico crítico: el contrato de transformación}

El contrato de transformación de la librería base expone una única operación, aplicable a
una muestra a la vez y sin estado. De forma coherente, el protocolo de evaluación
transformaba el conjunto de datos completo \emph{antes} de partirlo en entrenamiento y
validación, lo cual es inofensivo mientras las transformaciones no tengan estado.

El preprocesamiento tabular invalida ese supuesto. El escalado necesita la media y la
desviación de las columnas; la imputación, la media o la moda; la codificación, el catálogo
de categorías. Todas esas cantidades deben estimarse \textbf{únicamente con el subconjunto
de entrenamiento}: estimarlas sobre el conjunto completo introduce información de la
validación en la transformación aplicada al entrenamiento, sesga el desempeño medido y
anula la validez de los experimentos.

Este era el riesgo técnico central del proyecto, y se resolvió en el diseño antes de
escribir cualquier transformación tabular. La Tabla~\ref{tab:alternativas-contrato}
compara las alternativas consideradas.

\begin{table}[H]
\centering
\small
\caption{Alternativas consideradas para extender el contrato de transformación}
\label{tab:alternativas-contrato}
\begin{tabularx}{\textwidth}{>{\ancho{0.8}}L>{\ancho{0.48}}C>{\ancho{1.72}}L}
\toprule
\textbf{Alternativa} & \textbf{Decisión} & \textbf{Razón} \\
\midrule
Añadir ajuste y transformación como métodos abstractos al contrato existente &
Rechazada &
Obligaría a modificar las siete transformaciones de imagen y rompería cualquier
transformación personalizada ya escrita contra el contrato publicado. \\
\addlinespace
Sustituir la operación por muestra por una operación por lote &
Rechazada &
Cambia una interfaz pública en uso por los ejemplos y por la suite de regresión, sin
ganancia funcional. \\
\addlinespace
Definir un contrato separado para lo tabular, sin relación con el existente &
Rechazada &
Obligaría a duplicar la composición de pipelines y los motores de búsqueda, lo que
contradice la estrategia de evolución sobre reinvención. \\
\addlinespace
Añadir ajuste y transformación como métodos \textbf{concretos}, con implementación por
defecto sin estado, más una subclase que exige el ajuste de forma explícita &
\textbf{Adoptada} &
Ninguna transformación existente requiere cambios: heredan un ajuste que no aprende nada y
una transformación por lote que reproduce el comportamiento anterior. Las
transformaciones con estado heredan de la subclase, que registra el ajuste y rechaza
transformar antes de ajustar. \\
\bottomrule
\end{tabularx}
\end{table}

La decisión se acompañó de un segundo hallazgo obtenido al leer los motores de búsqueda:
cada candidato se construye tomando instancias del espacio de búsqueda, de modo que
\textbf{la misma instancia de transformación aparece en muchos candidatos}. Con
transformaciones sin estado eso es inofensivo, pero en cuanto el ajuste almacene estado en
el objeto, un candidato contaminaría a los siguientes, y bajo el backend de paralelización
habría además una condición de carrera entre procesos. La solución fue clonar el pipeline
dentro del protocolo de evaluación, por la misma razón por la que ya se clonaba el modelo
base.

El protocolo de evaluación resultante garantiza cuatro propiedades, cada una con su prueba
asociada: el pipeline se ajusta solo con el subconjunto de entrenamiento; el pipeline se
clona antes de ajustarse; el modelo base se clona; y todos los candidatos comparten la
misma partición, para que la comparación sea justa. Las dos primeras pruebas están escritas
para \emph{fallar} si alguien revierte el cambio: la de fuga usa un conjunto de datos donde
entrenamiento y validación tienen escalas deliberadamente distintas y verifica que la
transformación ajustada vio solo las filas de entrenamiento; la de aislamiento evalúa dos
candidatos que comparten una instancia con estado y comprueba que el resultado coincide con
evaluarlos por separado.

\subsubsection{Spikes de viabilidad}

Antes de construir el recomendador y el extractor de meta-características se ejecutaron dos
experimentos acotados para medir supuestos sobre los que descansaba el resto del proyecto.

\textbf{Spike A: inyección de ensayos iniciales.} Verificó que las configuraciones
inyectadas se consumen como primeros ensayos del estudio y en orden; que cuentan dentro del
presupuesto de ensayos; que la secuencia es reproducible con la misma semilla; y que el
podador configurado no interviene cuando el objetivo no reporta valores intermedios. También
documentó dos comportamientos peligrosos en direcciones opuestas: un índice fuera de rango
aborta el estudio completo, mientras que una clave desconocida se ignora en silencio. Ambos
obligan a validar cada configuración recomendada contra la versión vigente del espacio de
búsqueda antes de inyectarla.

El resultado más relevante del spike fue, sin embargo, otro. Comparando cinco formas de
arranque sobre un problema controlado, con diez semillas y dos geometrías distintas de la
función objetivo, se observó que el mecanismo funciona cuando las recomendaciones son
buenas —el objetivo se alcanza en el primer ensayo—, pero que \textbf{recomendaciones
plausibles y equivocadas resultan peores que no recomendar nada, e incluso peores que
recomendar al azar}. La explicación es el mecanismo del muestreador: varios puntos
agrupados con puntajes altos forman una región que el modelo interpreta como buena, y la
búsqueda se concentra ahí en lugar de explorar. Este resultado convierte en medible el
riesgo principal del proyecto y eleva la importancia de la ablación con un recomendador
aleatorio, ya prevista como análisis de control.

\textbf{Spike B: costo de la extracción de meta-características.} Midió el tiempo de
extracción por grupo de meta-características sobre tamaños crecientes. Los grupos general,
estadístico y teórico-informativo cuestan menos de tres segundos en conjunto incluso en el
techo de 50\,000 instancias. El grupo de complejidad, en cambio, escala con un exponente
medido de 1{,}91 —comportamiento cuadrático, coherente con medidas basadas en distancias— y
su costo extrapolado para un solo conjunto de datos en ese techo es de unos 22 minutos.

El spike identificó además que el grupo de complejidad no admite valores faltantes: lanza
una excepción sin mensaje. Imputar antes de extraer resuelve el fallo por completo. Este es
el riesgo de heterogeneidad de los conjuntos de datos públicos, anticipado en el
anteproyecto, materializado y con mitigación verificada.

Ambos spikes produjeron decisiones de configuración que se documentaron en el pre-registro
antes de que existiera cualquier resultado experimental: la fijación explícita del número de
ensayos iniciales del muestreador, y el submuestreo del grupo de complejidad a un tamaño fijo
igual para todos los conjuntos de datos.

\subsubsection{Núcleo agnóstico a la tarea}

La interfaz de tarea se implementó como una estructura de datos inmutable que agrupa las
cuatro responsabilidades identificadas por la auditoría: el conjunto de métricas válidas, la
dirección de optimización, la normalización de predicciones y la obtención de puntajes
continuos. Se descartó una jerarquía de clases porque una tarea no tiene comportamiento
propio que extender: es una colección de decisiones. La estructura valida su propia
coherencia al construirse y es serializable, requisito para que pueda viajar a los procesos
de trabajo de la evaluación paralela.

La tarea se selecciona en el constructor de los motores de búsqueda y del ejecutor paralelo,
nunca en la firma de la operación de búsqueda, de modo que ninguna interfaz pública en uso
cambió. Su valor por defecto es la detección de anomalías, con lo que el comportamiento del
caso de uso original se preserva de forma exacta. La métrica objetivo pasó a ser opcional: si
no se indica, se usa la métrica primaria de la tarea, lo que permite construir el motor para
clasificación y obtener F1 macro sin indicar nada más.

La Tabla~\ref{tab:tareas} contrasta las dos tareas implementadas.

\begin{table}[H]
\centering
\small
\caption{Comportamiento dependiente de la tarea}
\label{tab:tareas}
\begin{tabularx}{\textwidth}{>{\ancho{0.5}}L>{\ancho{1.25}}L>{\ancho{1.25}}L}
\toprule
\textbf{Aspecto} & \textbf{Detección de anomalías} & \textbf{Clasificación} \\
\midrule
Métricas &
AUC-ROC, F1, precisión, exhaustividad, tasa de falsos positivos &
F1 macro, exactitud, exactitud balanceada, precisión macro, exhaustividad macro \\
Métrica primaria & AUC-ROC & F1 macro \\
Se minimiza & Tasa de falsos positivos & Ninguna \\
Normalización de predicciones &
Se normalizan a dos clases, admitiendo la convención de signo de scikit-learn &
Se usan las etiquetas de clase tal cual, sin reinterpretar \\
Puntajes continuos &
Probabilidad de la clase positiva, o función de decisión con signo invertido &
No se usan; ninguna métrica macro se calcula sobre un ordenamiento \\
\bottomrule
\end{tabularx}
\end{table}

\subsubsection{Soporte de datos tabulares}

El cargador tabular admite un archivo de valores separados por comas o una tabla en memoria,
separa la columna objetivo, codifica las etiquetas a enteros conservando los valores
originales y reporta qué columnas son numéricas y cuáles categóricas. Para etiquetas que ya
son binarias la codificación es la identidad, con lo que se preserva la convención de la
detección de anomalías.

La decisión de diseño que gobierna el resto del soporte tabular es que el cargador devuelve
una \textbf{tabla con tipos} y no una matriz numérica. La razón no es de comodidad: las
transformaciones tabulares necesitan saber qué columnas son categóricas para codificarlas, y
si esa información se pasara como parámetro —por ejemplo, una lista de posiciones de
columna— el espacio de búsqueda dependería del conjunto de datos y las configuraciones
dejarían de ser transferibles entre conjuntos, que es precisamente lo que el arranque en
caliente necesita. Conservando los tipos, cada transformación detecta por sí misma sobre qué
columnas actuar.

Se implementaron tres transformaciones con estado: imputación de faltantes, escalado
numérico con tres métodos y codificación de categóricas con dos métodos. Todas heredan de
una base común que exige una tabla, recuerda las columnas del ajuste, valida que estén
presentes al transformar y resuelve el contrato por muestra. Ninguna recibe nombres ni
posiciones de columnas, propiedad que una prueba parametrizada vigila de forma permanente.

El punto más delicado fue la codificación. Las categorías se aprenden solo del subconjunto
de entrenamiento, de modo que una categoría que aparezca únicamente en validación no puede
crear una columna nueva: en codificación indicadora produce ceros en todas las columnas de
ese grupo, y en codificación ordinal recibe un código reservado. Es lo que garantiza que la
matriz de validación tenga el mismo ancho que la de entrenamiento. La forma directa de
resolver el problema —aplicar la codificación por separado a cada subconjunto— rompe esa
garantía en silencio y entrega al modelo matrices de anchos distintos.

Para que una tabla pudiera recorrer el protocolo de evaluación, dos funciones auxiliares
aprendieron a tratarla: la selección de subconjuntos usa indexación posicional, porque
recorrer una tabla con indexación ordinaria devolvería sus columnas en lugar de sus filas, y
la conversión a matriz de características la convierte de una sola vez. Ambas detectan la
tabla \emph{por capacidades y no por tipo}, de modo que los módulos del núcleo no importan
la librería de manipulación de tablas: el soporte tabular es opcional y el camino de
imágenes no debe pagar ese costo de importación.

\subsubsection{Criterios de inclusión del corpus}

La selección del corpus parte de la colección curada OpenML-CC18 \cite{openml-cc18}, de 72
conjuntos de datos, en lugar de filtrar el repositorio completo de OpenML. La razón es de
reproducibilidad: el repositorio completo contiene versiones duplicadas, metadatos
incompletos y cargas defectuosas, de modo que filtrarlo por tamaño produciría una lista que
habría que depurar a mano sin criterio replicable. La colección curada tiene metadatos
verificados e identidad estable.

Aplicando los criterios de inclusión sobre esa colección quedan 52 conjuntos de datos
candidatos. Los descartes se reparten en seis por número de instancias fuera del rango
admitido, quince por exceder las 100 características predictoras y cinco por número de
clases fuera del rango. La generación de la lista está automatizada en un script que
consulta la interfaz de OpenML, de modo que puede regenerarse y auditarse.

La revisión de la lista resultante produjo un hallazgo con implicaciones sobre la validez
del protocolo experimental. El corpus candidato contiene \textbf{familias de conjuntos de
datos que comparten origen}, y en un caso comparten las propias instancias: cuatro de los
candidatos son cuatro descripciones distintas del mismo conjunto de 2\,000 imágenes de
dígitos manuscritos, cada una con un conjunto de características diferente. Si el corpus
incluyera varias de ellas y la validación por exclusión dejara una fuera, el recomendador
seguiría teniendo ese mismo conjunto de datos en su base de entrenamiento bajo otra
representación: el supuesto de independencia entre pliegues dejaría de cumplirse y la
ventaja medida del arranque en caliente saldría inflada por una vía difícil de detectar en
los resultados finales.

En consecuencia se añadió una regla de selección de a lo sumo un conjunto de datos por
familia de origen, con dos niveles: obligatorio cuando las instancias son las mismas, por
tratarse de un problema de validez, y recomendable cuando solo se comparte la fuente o el
esquema de características, por tratarse de una cuestión de diversidad. Aplicada a los 52
candidatos, la regla deja 42 familias distintas, margen suficiente para seleccionar los 15 a
25 conjuntos de datos comprometidos.

\subsubsection{Estado de los entregables}

La Tabla~\ref{tab:estado-entregables} resume el estado de los entregables a la fecha.

\begin{table}[H]
\centering
\small
\caption{Estado de los entregables a la fecha de redacción}
\label{tab:estado-entregables}
\begin{tabularx}{\textwidth}{>{\ancho{0.4}}C>{\ancho{1.3}}L>{\ancho{0.55}}C>{\ancho{1.75}}L}
\toprule
\textbf{\#} & \textbf{Entregable} & \textbf{Estado} & \textbf{Evidencia} \\
\midrule
RE1 & Núcleo agnóstico a la tarea & Completo &
Interfaz de tarea con dos tareas implementadas; protocolo de evaluación con ajuste solo
sobre entrenamiento y aislamiento entre candidatos; suite de regresión en verde. \\
\addlinespace
RE2 & Soporte de datos tabulares & Parcial &
Cargador y tres transformaciones con estado, con flujo completo verificado. Pendiente la
selección de características y el espacio de búsqueda canónico. \\
\addlinespace
RE3 & Extractor de meta-características & No iniciado &
Viabilidad y costo medidos en el spike B; decisiones de configuración registradas. \\
\addlinespace
RE4 & Base de meta-conocimiento & Parcial &
Criterios de inclusión definidos y lista de 52 candidatos generada de forma reproducible. \\
\addlinespace
RE5 & Recomendador meta-aprendido & No iniciado &
Mecanismo de inyección verificado en el spike A. \\
\addlinespace
RE6 & Experimentos y reportes & No iniciado &
Protocolo pre-registrado; decisiones de medición fijadas. \\
\addlinespace
RE7 & Documentación y código & Parcial &
Guía técnica, documentos de decisión de arquitectura, informes de spikes y ejemplo
ejecutable del flujo tabular. \\
\bottomrule
\end{tabularx}
\end{table}

\subsection{Análisis de los resultados obtenidos}

\subsubsection{Criterios de éxito de ingeniería: estado verificado}

Los criterios de éxito de ingeniería son independientes del resultado de la hipótesis
experimental y ya pueden evaluarse. La Tabla~\ref{tab:criterios-ingenieria} presenta su
estado medido a la fecha de redacción.

\begin{table}[H]
\centering
\small
\caption{Estado de los criterios de éxito de ingeniería}
\label{tab:criterios-ingenieria}
\begin{tabularx}{\textwidth}{>{\ancho{1.18}}L>{\ancho{1.0}}L>{\ancho{0.57}}C>{\ancho{1.25}}L}
\toprule
\textbf{Criterio} & \textbf{Medición} & \textbf{Estado} & \textbf{Observación} \\
\midrule
La suite de regresión de la librería base se mantiene en verde tras la refactorización &
Las 134 pruebas migradas pasan sin haber sido editadas &
Cumplido &
Es la evidencia central de que la refactorización no alteró el caso de uso original: las
pruebas afirman valores numéricos concretos de las métricas. \\
\addlinespace
Los módulos nuevos alcanzan al menos un 70\,\% de cobertura &
Cobertura global de 99\,\% sobre 993 instrucciones; los módulos de tarea, cargador tabular
y transformaciones tabulares al 100\,\% &
Cumplido &
La suite pasó de 134 a 287 pruebas y la cobertura global de 97\,\% a 99\,\%. \\
\addlinespace
Los resultados son reproducibles: la misma semilla produce el mismo pipeline ganador &
Verificado por las pruebas de reproducibilidad de la librería base y por la fijación
explícita de los parámetros del muestreador &
Parcial &
Verificado a nivel de componente. La verificación de punta a punta sobre el experimento
completo queda pendiente de su ejecución. \\
\addlinespace
El flujo completo se mantiene expresable en no más de 20 líneas de código &
El flujo tabular completo usa 13 instrucciones &
Cumplido &
La verificación está automatizada: una prueba cuenta las instrucciones del ejemplo y lo
ejecuta, de modo que el criterio se comprueba en cada corrida de la suite en lugar de
afirmarse. \\
\bottomrule
\end{tabularx}
\end{table}

Sobre el último criterio conviene precisar el método de conteo, porque admite más de una
lectura. Se cuentan \emph{instrucciones} y no líneas físicas: el límite de longitud de línea
que impone el estilo del proyecto obliga a repartir una sola instrucción en varias líneas,
de modo que contar líneas mediría el estilo de formato en lugar de la complejidad del flujo.

\guia{Analice los resultados obtenidos.}

\guia{Incluya:}
\begin{itemize}
  \item \guia{Evaluaciones}
  \item \guia{Métricas}
  \item \guia{Comparaciones}
  \item \guia{Validación de objetivos}
\end{itemize}

\guia{Apóyese con tablas y gráficos.}

\clearpage
%------------------------------------------------------------------------------
\section{Conclusiones y recomendaciones}
\label{sec:conclusiones}

\subsection{Conclusiones}

\guia{En esta sección se indica si se cumplieron los objetivos. Es importante
destacar los principales logros del proyecto.}

\subsection{Recomendaciones}

\guia{Aquí se sugieren mejoras o trabajos futuros. Se debe incluir aspectos no
desarrollados pero relevantes.}

\clearpage
%------------------------------------------------------------------------------
\section{Apéndices y anexos}
\label{sec:apendices}

\subsection*{Apéndice 1: Implicaciones ambientales y de desarrollo sostenible}
\addcontentsline{toc}{subsection}{Apéndice 1: Implicaciones ambientales y de desarrollo sostenible}

\guia{Este apéndice se debe de realizar de manera obligatoria y debe de incluir
una reflexión objetiva y con amplia proyección respecto del entorno sobre las
implicaciones e impactos actuales y potenciales relacionados con aspectos
sociales, económicos, ambientales, legales, de seguridad y salud que conlleva
su proyecto de ingeniería para la sociedad en general, sobre todo para la
práctica de la ingeniería y el ejercicio profesional del ingeniero en
computadores. Recordar que se debe de entregar una copia digital de este
apéndice en un documento por separado, el cual debe de tener la portada
original del informe y se entrega en formato PDF.}

\guia{Para el resto de los apéndices, se debe incluir el material
complementario utilizado, como por ejemplo documentos, instructivos,
instrumentos de encuestas, programas, algoritmos, entre otros.}

\clearpage
%------------------------------------------------- REFERENCIAS BIBLIOGRÁFICAS -
\section{Referencias bibliográficas}
\label{sec:referencias}

\guia{Esta sección incluye la lista de materiales (libros, revistas, sitios
web, etc.) y otros medios de recolección de información (congresos,
entrevistas, tesis, etc.) que fueron consultados. A lo largo de todo el
documento, las citaciones de referencias en el texto deberán usar el estándar
IEEE de manera rigurosa.}

\guia{La lista se genera automáticamente a partir de \texttt{referencias.bib}:
agregue ahí cada fuente y cítela en el texto con \cmd{cite\{clave\}}. Solo
aparecen aquí las entradas efectivamente citadas.}

\printbibliography[heading=none]

%==============================================================================
\end{document}
%==============================================================================
